% GENERATED FILE. Edit canonical lessons, then run npm run generate:books.
% canonical-source-hash: fnv1a64:68759e5ce682ec74
% canonical-lessons: TE-C160-sampadincu, TE-C160-vivarincu, TE-C160-anuvadincu, TE-C160-gurtupettuko, TE-C160-oppuko

\chapter{To earn, To explain, To translate, To remember, To agree}
\label{ch:te-to-earn-to-explain-to-translate-to-remember-to-agree}
\hlchaptermodality{160}

\begin{chapteropening}
\textbf{By the end of this chapter:} \emph{I can say to earn, to explain, to translate, to remember and to agree.}

Complete the last lesson of chapter 160: I can say to earn, to explain, to translate, to remember and to agree.
\end{chapteropening}

\section[\texorpdfstring{\te{సంపాదించు} (sampādiñcu)}{సంపాదించు (sampādiñcu)}]{\te{సంపాదించు} (sampādiñcu) — to earn}

\label{lesson:TE-C160-sampadincu}



\begin{quote}
Before the new one: say the Telugu for to throw away, then the Telugu for to lose.
\end{quote}

\subsection*{You'll want to know: \te{సంపాదించు}}
\textbf{\te{సంపాదించు}} — \emph{sampādiñcu} — \textquotedblleft{}to earn\textquotedblright{}.

\begin{grammarlens}[title={What you've built}]
One more word for this chapter.
\end{grammarlens}

\subsection*{Your turn}
\begin{itemize}
\raggedright
  \item \emph{Say it:} \emph{sampādiñcu}
  \item \emph{Say it:} \emph{sampādiñcu}, once more
  \item \emph{Say it:} say \emph{sampādiñcu}
  \item \emph{Recall:} say the Telugu for to throw away, then the Telugu for to lose, then say \emph{sampādiñcu} again
\end{itemize}

\begin{tcolorbox}[breakable,colback=teal!4,colframe=teal!35!black,title={Before you move on}]
What does \te{సంపాదించు} mean? (\textquotedblleft{}To earn\textquotedblright{}.) Say it once more.
\end{tcolorbox}
\section[\texorpdfstring{\te{వివరించు} (vivariñcu)}{వివరించు (vivariñcu)}]{\te{వివరించు} (vivariñcu) — to explain}

\label{lesson:TE-C160-vivarincu}



\begin{quote}
Before the new one: say the Telugu for to lose, then the Telugu for to earn.
\end{quote}

\subsection*{You'll want to know: \te{వివరించు}}
\textbf{\te{వివరించు}} — \emph{vivariñcu} — \textquotedblleft{}to explain\textquotedblright{}.

\begin{grammarlens}[title={What you've built}]
One more word for this chapter.
\end{grammarlens}

\subsection*{Your turn}
\begin{itemize}
\raggedright
  \item \emph{Say it:} \emph{vivariñcu}
  \item \emph{Say it:} \emph{vivariñcu}, once more
  \item \emph{Say it:} say \emph{vivariñcu}
  \item \emph{Recall:} say the Telugu for to lose, then the Telugu for to earn, then say \emph{vivariñcu} again
\end{itemize}

\begin{tcolorbox}[breakable,colback=teal!4,colframe=teal!35!black,title={Before you move on}]
What does \te{వివరించు} mean? (\textquotedblleft{}To explain\textquotedblright{}.) Say it once more.
\end{tcolorbox}
\section[\texorpdfstring{\te{అనువదించు} (anuvadiñcu)}{అనువదించు (anuvadiñcu)}]{\te{అనువదించు} (anuvadiñcu) — to translate}

\label{lesson:TE-C160-anuvadincu}



\begin{quote}
Before the new one: say the Telugu for to earn, then the Telugu for to explain.
\end{quote}

\subsection*{You'll want to know: \te{అనువదించు}}
\textbf{\te{అనువదించు}} — \emph{anuvadiñcu} — \textquotedblleft{}to translate\textquotedblright{}.

\begin{grammarlens}[title={What you've built}]
One more word for this chapter.
\end{grammarlens}

\subsection*{Your turn}
\begin{itemize}
\raggedright
  \item \emph{Say it:} \emph{anuvadiñcu}
  \item \emph{Say it:} \emph{anuvadiñcu}, once more
  \item \emph{Say it:} say \emph{anuvadiñcu}
  \item \emph{Recall:} say the Telugu for to earn, then the Telugu for to explain, then say \emph{anuvadiñcu} again
\end{itemize}

\begin{tcolorbox}[breakable,colback=teal!4,colframe=teal!35!black,title={Before you move on}]
What does \te{అనువదించు} mean? (\textquotedblleft{}To translate\textquotedblright{}.) Say it once more.
\end{tcolorbox}
\section[\texorpdfstring{\te{గుర్తుపెట్టుకో} (gurtupeṭṭukō)}{గుర్తుపెట్టుకో (gurtupeṭṭukō)}]{\te{గుర్తుపెట్టుకో} (gurtupeṭṭukō) — to remember, to keep in mind}

\label{lesson:TE-C160-gurtupettuko}



\begin{quote}
Before the new one: say the Telugu for to explain, then the Telugu for to translate.
\end{quote}

\subsection*{You'll want to know: \te{గుర్తుపెట్టుకో}}
\textbf{\te{గుర్తుపెట్టుకో}} — \emph{gurtupeṭṭukō} — \textquotedblleft{}to remember, to keep in mind\textquotedblright{}.

\begin{grammarlens}[title={What you've built}]
One more word for this chapter.
\end{grammarlens}

\subsection*{Your turn}
\begin{itemize}
\raggedright
  \item \emph{Say it:} \emph{gurtupeṭṭukō}
  \item \emph{Say it:} \emph{gurtupeṭṭukō}, once more
  \item \emph{Say it:} say \emph{gurtupeṭṭukō}
  \item \emph{Recall:} say the Telugu for to explain, then the Telugu for to translate, then say \emph{gurtupeṭṭukō} again
\end{itemize}

\begin{tcolorbox}[breakable,colback=teal!4,colframe=teal!35!black,title={Before you move on}]
What does \te{గుర్తుపెట్టుకో} mean? (\textquotedblleft{}To remember, to keep in mind\textquotedblright{}.) Say it once more.
\end{tcolorbox}
\section[\texorpdfstring{\te{ఒప్పుకో} (oppukō)}{ఒప్పుకో (oppukō)}]{\te{ఒప్పుకో} (oppukō) — to agree, to accept}

\label{lesson:TE-C160-oppuko}



\begin{quote}
Before the new one: say the Telugu for to translate, then the Telugu for to remember.
\end{quote}

\subsection*{You'll want to know: \te{ఒప్పుకో}}
\textbf{\te{ఒప్పుకో}} — \emph{oppukō} — \textquotedblleft{}to agree, to accept\textquotedblright{}.

\begin{grammarlens}[title={What you've built}]
One more word for this chapter.
\end{grammarlens}

\subsection*{Your turn}
\begin{itemize}
\raggedright
  \item \emph{Say it:} \emph{oppukō}
  \item \emph{Say it:} \emph{oppukō}, once more
  \item \emph{Say it:} say \emph{oppukō}
  \item \emph{Recall:} say the Telugu for to translate, then the Telugu for to remember, then say \emph{oppukō} again
\end{itemize}

\begin{tcolorbox}[breakable,colback=teal!4,colframe=teal!35!black,title={Before you move on}]
What does \te{ఒప్పుకో} mean? (\textquotedblleft{}To agree, to accept\textquotedblright{}.) Say it once more.
\end{tcolorbox}
